
% --- SLIDE 6: REST VS GRPC + LATENCY ---
\begin{frame}{Protocols and Latency Budgets}

    \begin{columns}[T]
        \begin{column}{0.46\textwidth}
            \small
            \textbf{REST/JSON vs gRPC}
            \begin{itemize}
                \item \textbf{REST + JSON:} human-readable, debuggable with \texttt{curl}, works with every
                      client. Verbose; parsing costs real time on large payloads.
                \item \textbf{gRPC + protobuf:} compact binary, schema-checked, streaming, faster for
                      high-volume internal traffic. Harder to inspect.
                \item \textbf{Pick REST first.} Move to gRPC when profiling says serialisation is your
                      bottleneck --- rarely true for a scikit-learn model.
            \end{itemize}
        \end{column}
        \begin{column}{0.5\textwidth}
            \begin{tcolorbox}[colback=blue!5, colframe=blue!50, title=\small Budget the whole request]
                \scriptsize
                \begin{tabular}{lr}
                    Network in/out        & 5--20\,ms \\
                    Feature lookup (DB / feature store) & 5--50\,ms \\
                    Preprocessing         & 1--5\,ms \\
                    \textbf{Model forward pass} & \textbf{1--10\,ms} \\
                    Post-processing, logging & 1--5\,ms \\
                \end{tabular}
                \vspace{0.4em}
                \\ The model is often the \emph{smallest} term. Optimise what you measured, not what you
                assumed.
            \end{tcolorbox}
            \vspace{0.4em}
            \small \textbf{Quote p95 and p99, never the mean.} The tail is what users experience and what
            times out upstream.
        \end{column}
    \end{columns}

    \vspace{0.5em}
    \begin{block}{Throughput vs latency}
        \small Batching requests on the server raises throughput and \emph{raises} latency. Pick which one
        the product actually needs, then set the batch window accordingly.
    \end{block}

\end{frame}
